% v2-figures/town-map.tex: the walled-town plan, with the real component beside
% every civic name. The analogy describes structure, not political centralism:
% the second town to the right is another authority with its own hall and its
% own wall, joined by a road that runs one way.
\begin{tikzpicture}[
  wall/.style={draw=navydark, line width=2.2pt, fill=tabstripe},
  district/.style={draw=navymid, line width=0.7pt, rounded corners=2pt, fill=white, font=\scriptsize, align=center, inner sep=3pt, text width=2.8cm},
  hall/.style={draw=accent, line width=1.4pt, rounded corners=2pt, fill=accent, text=white, font=\bfseries\scriptsize, align=center, inner sep=4pt},
  gate/.style={draw=navydark, line width=1pt, fill=gold, font=\tiny\bfseries, align=center, inner sep=2pt},
  tower/.style={draw=navydark, line width=1pt, fill=navydark, text=white, font=\tiny\bfseries, align=center, inner sep=2.5pt, circle},
  law/.style={draw=classmark, line width=0.8pt, fill=white, text=classmark, font=\scriptsize\bfseries, align=center, inner sep=4pt, text width=10.4cm},
  road/.style={draw=steelgray, line width=3pt, double=white, double distance=1pt},
  civic/.style={font=\tiny\scshape\color{gold}, inner sep=1pt},
]
% The laws above the town.
\node[law] (law) at (0,6.6) {The laws above the town: the vocation and C1--C10 \mbox{(MISSION.md)}, enforced in the schema, pinned by 191 checks};
\draw[flowfaint] (law.south) -- (0,5.5);
% The wall (schema constraints).
\draw[wall] (-5.4,-4.8) rectangle (5.4,5.0);
\node[civic, anchor=south west, align=left] at (-5.12,5.18) {the wall: triggers, checks,\\unique indexes};
% Districts.
\node[hall, minimum width=3.2cm] (hall) at (0,1.2) {town hall record\\IdentityToken + TokenSignature};
\node[district] (work) at (-3.3,3.3) {workshops\\epochs, permissions,\\recovery, duress};
\node[district] (council) at (3.3,3.3) {council chamber\\{\tiny AgencyAlgorithmAuth,}\\{\tiny IssuerDiscretionPolicy, quotas}};
\node[district] (arch) at (-3.3,-1.5) {archives\\{\tiny append-only event logs,}\\{\tiny RetentionPolicy}};
\node[district] (reg) at (3.3,-1.5) {registry office\\Athena views,\\the signed registry,\\the trust list};
\node[district, text width=3.4cm] (court) at (0,-3.5) {court records\\{\tiny AuditAccessLog, DuressEvent,}\\{\tiny LifecycleArchiveCheckpoint}};
% Gates on the wall.
\node[gate] (g1) at (0,5.0) {issuance gate\\\texttt{uc1\_issue}, signing};
\node[gate] (g2) at (5.4,1.2) {verification gate\\\texttt{/api/v1/verify}};
\node[gate] (g3) at (0,-4.8) {exchange gate\\\texttt{/api/v1/exchange}};
\node[gate, anchor=west] (g4) at (-5.4,1.2) {login gate\\auth broker, PKCE};
% Watchtowers at the corners.
\node[tower] (t1) at (-5.4,5.0) {W};
\node[tower] (t2) at (5.4,5.0) {W};
\node[tower] (t3) at (-5.4,-4.8) {W};
\node[tower] (t4) at (5.4,-4.8) {W};
\node[civic, anchor=north west] at (-5.4,-5.2) {W: watchtowers, the transparency logs};
% Watchers outside the wall.
\node[extbox, minimum width=2.6cm] (wit) at (-9.4,-3.4) {independent witnesses\\[-1pt]{\tiny cosign heads, gossip,}\\{\tiny refuse a fork}};
\draw[sees] (wit.east) -- (t3.west);
\draw[sees] (wit.north east) -- (t1.south west);
\node[extbox, minimum width=2.6cm] (rp) at (8.5,1.2) {relying party\\[-1pt]{\tiny verifies offline with}\\{\tiny the detached verifier}};
\draw[flow] (g2.east) -- (rp.west);
% The second town: another authority, joined by a one-way road.
\draw[wall, fill=white] (7.2,-4.6) rectangle (10.8,-2.2);
\node[hall, minimum width=2.4cm, font=\bfseries\tiny] at (9.0,-3.4) {another authority\\its own hall, its own key};
\draw[road] (5.4,-3.4) -- (7.2,-3.4);
\node[civic, anchor=south, align=center] at (6.3,-3.1) {road:\\attestation};
\node[font=\tiny\itshape\color{steelgray}, anchor=north, align=center] at (6.3,-3.7) {one way,\\one context};
% The holder outside.
\node[humanbox, minimum width=2.4cm] (holder) at (-9.4,1.2) {holder\\[-1pt]{\tiny wallet file, presentation}};
\draw[flow] (holder.east) -- (g4.west);
\node[font=\tiny\itshape\color{steelgray}, anchor=north, text width=3.0cm, align=center] at (-9.4,0.3) {enters by possession; a duress code looks like consent to the coercer};
\end{tikzpicture}
